% ===== Block B09: 原文 p.73–82 | §4.1 共形块层 + §4.2 局部自由性 I（光滑情形） | 译手 tr-ch3b =====

% [原文 p.73]
\subsection{共形块层}

{\emergencystretch=2.8pt 本节我们将定义带形式坐标（formal coordinate）的稳定 $N$-点曲线（stable
$N$-pointed curve）族 $\mathcal{F}$ 的共形块层（sheaf of conformal blocks）
$\cV_\lambda^\dagger(\mathcal{F})$。\par}

设 $\mathcal{F}=(\pi:\mathcal{C}\to\mathcal{B};s_1,\dots,s_N;
\eta_1,\dots,\eta_N)$ 是亏格为 $g$ 的带形式坐标的 $N$-点稳定曲线族。我们假定
$\mathcal{B}$ 是有限维复流形（complex manifold），且族的每个纤维（fibre）都满足
§1.1 中的条件 (Q)，但我们\emph{不}假定该族连通。本节的主要目的是定义附着于
该族的真空层（sheaf of vacua）$\cV_\lambda^\dagger(\mathcal{F})$，并证明它是
凝聚（coherent）$\cO_\mathcal{B}$-模。

\begin{definition}[4.1]
$\mathcal{B}$ 上的仿射李代数层（sheaf of affine Lie algebras）
$\hat{\mathfrak{g}}_N(\mathcal{B})$ 是 $\cO_\mathcal{B}$-模层
\[
  \hat{\mathfrak{g}}_N(\mathcal{B})
  = \mathfrak{g}\otimes_\C\Bigl(\bigoplus_{j=1}^N
    \cO_\mathcal{B}((\xi_j))\Bigr)\oplus\cO_\mathcal{B}\cdot c ,
\]
其交换关系如下，它是 $\cO_\mathcal{B}$-双线性（bilinear）的：
\[
  \begin{aligned}
  &[(X_1\otimes f_1,\dots,X_N\otimes f_N),
    (Y_1\otimes g_1,\dots,Y_N\otimes g_N)]\\
  &\quad = ([X_1,Y_1]\otimes(f_1g_1),\dots,[X_N,Y_N]\otimes(f_Ng_N))\\
  &\qquad \oplus\Bigl(\sum_{j=1}^N(X_j,Y_j)
    \operatorname{Res}_{\xi_j=0}(g_jdf_j)\Bigr)\cdot c ,
  \end{aligned}
\]
其中 $c$ 属于中心（center），且
\[
  X_j,Y_j\in\mathfrak{g},\qquad f_j,g_j\in\cO_\mathcal{B}((\xi_j)) .
\]
置
\begin{equation}
  \hat{\mathfrak{g}}(\mathcal{F})
  = \mathfrak{g}\otimes_\C\pi_*(\cO_\mathcal{C}(*S)), \tag{4.1}
\end{equation}
其中我们定义
\[
  \begin{aligned}
  S &= \sum_{j=1}^N s_j(\mathcal{B}),\\
  \pi_*(\cO_\mathcal{C}(*S)) &= \varinjlim_k\,\pi_*(\cO_\mathcal{C}(kS)) .
  \end{aligned}
\]
利用形式邻域（formal neighbourhood）$\eta_j$，1.47 中定义的同态有层化版本
\[
  \tilde{t} : \pi_*(\cO_\mathcal{B}(*S))\to\bigoplus_{j=1}^N
  \cO_\mathcal{B}((\xi_j)) ,
\]
并且我们可以把 $\hat{\mathfrak{g}}(\mathcal{F})$ 看作
$\hat{\mathfrak{g}}_N(\mathcal{B})$ 的李子代数（Lie subalgebra）。
\end{definition}
% [Tutor Note: (4.1) 的 $\pi_*(\cO_\mathcal{C}(*S))$ 是「在 $S$ 上允许任意阶极点
%  的、沿纤维方向的函数层」被推到 $\mathcal{B}$ 上的直像：在每一点 $b\in\mathcal{B}$
%  附近，它的截面是 $\mathcal{C}$ 上在 $S$ 外全纯、在 $S$ 附近只有极点的函数。
%  $\varinjlim_k$ 只是把这些极点阶数不设上限地并起来（$*S=\cup_k kS$）。
%  下一行 $\tilde t$ 的定义域原文印成 $\pi_*(\cO_\mathcal{B}(*S))$，按 (4.1) 应为
%  $\pi_*(\cO_\mathcal{C}(*S))$，见回报 ERRATA。原文 p.73]
% [Tutor Note: $\tilde t$ 就是 1.47 中「把函数在每个标点 $s_j$ 的形式邻域里按
%  $\xi_j$ 作洛朗展开」的同态，只不过现在族在动、$\xi_j$ 是沿纤维的形式坐标，
%  于是它成为 $\cO_\mathcal{B}$-模层之间的同态。$\tilde t$ 是单射（在 $\mathcal{B}$
%  的每个点、每条纤维上，非零亚纯函数不可能在所有标点的形式邻域里都展开为 $0$），
%  这一点在下文的「李子代数」陈述里暗含：$\hat{\mathfrak{g}}(\mathcal{F})$ 通过
%  $\tilde t$ 嵌入 $\oplus_j\cO_\mathcal{B}((\xi_j))$。原文 p.73]

% [原文 p.74]
固定一个非负整数 $\ell$。对任意
$\vec{\lambda}=(\lambda_1,\dots,\lambda_N)\in(P_\ell)^N$，我们定义
\begin{equation}
  \cH_{\vec{\lambda}}(\mathcal{B})
  = \cO_\mathcal{B}\otimes_\C\cH_{\vec{\lambda}}, \tag{4.2}
\end{equation}
\begin{equation}
  \cH_{\vec{\lambda}}^\dagger(\mathcal{B})
  = \underline{\operatorname{Hom}}_{\cO_\mathcal{B}}
    (\cH_{\vec{\lambda}}(\mathcal{B}),\cO_\mathcal{B}) . \tag{4.3}
\end{equation}
配对 (2.26) 诱导出 $\cO_\mathcal{B}$-双线性配对
\begin{equation}
  \langle\ |\ \rangle : \cH_{\vec{\lambda}}^\dagger(\mathcal{B})
  \times\cH_{\vec{\lambda}}(\mathcal{B})\to\cO_\mathcal{B}, \tag{4.4}
\end{equation}
它关于下面引入的滤过（filtration）是完美（perfect）的。仿射李代数层
$\hat{\mathfrak{g}}_N(\mathcal{B})$ 通过
\begin{equation}
  \begin{aligned}
  &\Bigl(\bigl(X_1\otimes\sum_{n\in\Z}a_n^{(1)}\xi_1^n\bigr),\dots,
    \bigl(X_N\otimes\sum_{n\in\Z}a_n^{(N)}\xi_N^n\bigr)\Bigr)
    (F\otimes|\Psi\rangle)\\
  &\quad = \sum_{j=1}^N\sum_{n\in\Z}(a_n^{(j)}F)
    \otimes\rho_j(X_j(n))|\Psi\rangle
  \end{aligned}
  \tag{4.5}
\end{equation}
作用在 $\cH_{\vec{\lambda}}(\mathcal{B})$ 与
$\cH_{\vec{\lambda}}^\dagger(\mathcal{B})$ 上。
$\hat{\mathfrak{g}}_N(\mathcal{B})$ 在
$\cH_{\vec{\lambda}}^\dagger(\mathcal{B})$ 上的作用就是
$\cH_{\vec{\lambda}}(\mathcal{B})$ 的对偶作用（dual action），即
\[
  \langle\Psi a|\Phi\rangle = \langle\Psi|a\Phi\rangle
  \qquad\text{对任意 } a\in\hat{\mathfrak{g}}_N .
\]
% [Tutor Note: (4.5) 的左端是把 $\hat{\mathfrak{g}}_N(\mathcal{B})$ 的元素写成
%  $N$ 个分量，第 $j$ 个分量 $X_j\otimes\sum_n a_n^{(j)}\xi_j^n$ 中
%  $a_n^{(j)}$ 是 $\mathcal{B}$ 上的函数（即 $\cO_\mathcal{B}$ 的截面），
%  所以右端的 $(a_n^{(j)}F)$ 是「函数乘函数」，把 $\cO_\mathcal{B}$-模结构
%  与 $\mathfrak{g}$ 的作用真正耦合起来。若 $N=1$、$\mathcal{B}$ 是一点，
%  这就退化为 §2 中 $\hat{\mathfrak{g}}$ 在 $V_\lambda$ 上的标准作用。原文 p.74]

\begin{definition}[4.2]
置
\[
  \begin{aligned}
  \cV_{\vec{\lambda}}(\mathcal{F})
  &= \cH_{\vec{\lambda}}(\mathcal{B})
     \big/\hat{\mathfrak{g}}(\mathcal{F})\cH_{\vec{\lambda}}(\mathcal{B}),\\
  \cV_{\vec{\lambda}}^\dagger(\mathcal{F})
  &= \underline{\operatorname{Hom}}_{\cO_\mathcal{B}}
     (\cV_{\vec{\lambda}}(\mathcal{F}),\cO_\mathcal{B}) .
  \end{aligned}
\]
它们是 $\mathcal{B}$ 上 $\cO_\mathcal{B}$-模层。层
$\cV_{\vec{\lambda}}^\dagger(\mathcal{F})$ 称为附着于族 $\mathcal{F}$ 的
\emph{共形块层}或\emph{真空层}，而 $\cV_{\vec{\lambda}}(\mathcal{F})$ 称为
\emph{余真空层}（sheaf of covacua）。注意我们有
\[
  \cV_{\vec{\lambda}}^\dagger(\mathcal{F})
  = \{\langle\Psi|\in\cH_{\vec{\lambda}}^\dagger(\mathcal{B})\mid
     \langle\Psi|a=0 \ \text{对任意 } a\in\hat{\mathfrak{g}}(\mathcal{F})\} .
\]
配对 (4.4) 诱导出非退化（non-degenerate）$\cO_\mathcal{B}$-双线性配对
\begin{equation}
  \langle\ |\ \rangle : \cV_{\vec{\lambda}}^\dagger(\mathcal{F})
  \times\cV_{\vec{\lambda}}(\mathcal{F})\to\cO_\mathcal{B} . \tag{4.6}
\end{equation}
\end{definition}
% [Tutor Note: 定义 4.2 里「商再取对偶」的顺序不能颠倒：
%  $\cV_{\vec{\lambda}}(\mathcal{F})=\cH_{\vec{\lambda}}(\mathcal{B})/
%  \hat{\mathfrak{g}}(\mathcal{F})\cH_{\vec{\lambda}}(\mathcal{B})$ 是余不变量
%  （covacua），它的对偶才是真空层。取对偶后「$\hat{\mathfrak{g}}(\mathcal{F})$
%  作用为零」就变成「$\langle\Psi|a=0$」，这正是下一段那个集合等式。
%  又因为 (4.4) 是完美配对而 (4.6) 只是非退化配对：取商之后可能残余挠元，
%  所以 (4.6) 一般不再是完美配对——这一点决定了 §4.2 里能不能「局部自由」。
%  原文 p.74]

\begin{lemma}[4.3]
对点 $s\in\mathcal{B}$ 置
\[
  \begin{aligned}
  \mathcal{F}_s &:= (\pi^{-1}(s);s_1(s),\dots,s_N(s);
    \eta_1|_{\pi^{-1}(s)},\dots,\eta_N|_{\pi^{-1}(s)}),\\
  \C_s &:= \cO_{\mathcal{B},s}/\mathfrak{m}_s ,
  \end{aligned}
\]
其中 $\mathfrak{m}_s$ 是茎（stalk）$\cO_{\mathcal{B},s}$ 的极大理想（maximal
ideal）。于是我们有如下典范同构（canonical isomorphism）：
\[
  \begin{aligned}
  \C_s\otimes_{\cO_\mathcal{B}}\cH_{\vec{\lambda}}(\mathcal{B})
  &\simeq \cH_{\vec{\lambda}},\\
  \C_s\otimes_{\cO_\mathcal{B}}\hat{\mathfrak{g}}_N(\mathcal{B})
  &\simeq \hat{\mathfrak{g}}_N,\\
  \hat{\mathfrak{g}}(\mathcal{F})\otimes_{\cO_\mathcal{B}}\C_s
  &\simeq \hat{\mathfrak{g}}(\mathcal{F}_s),\\
  \cV_{\vec{\lambda}}(\mathcal{F})\otimes_{\cO_\mathcal{B}}\C_s
  &\simeq \cV_{\vec{\lambda}}(\mathcal{F}_s) .
  \end{aligned}
\]
更一般地，对全纯映射（holomorphic mapping）$f:Y\to\mathcal{B}$，记
$\mathcal{F}_Y$ 为族 $\mathcal{F}$ 沿态射（morphism）$f$ 的拉回
（pull-back）。于是我们有如下典范

% [原文 p.75]
同构：
\[
  \begin{aligned}
  \cO_Y\otimes_{\cO_\mathcal{B}}\cH_{\vec{\lambda}}(\mathcal{B})
  &\simeq \cH_{\vec{\lambda}}(Y),\\
  \hat{\mathfrak{g}}_N(\mathcal{B})\otimes_{\cO_\mathcal{B}}\cO_Y
  &\simeq \hat{\mathfrak{g}}_N(Y),\\
  \hat{\mathfrak{g}}(\mathcal{F})\otimes_{\cO_\mathcal{B}}\cO_Y
  &\simeq \hat{\mathfrak{g}}(\mathcal{F}_Y),\\
  \cO_Y\otimes_{\cO_\mathcal{B}}\cV_{\vec{\lambda}}(\mathcal{F})
  &\simeq \cV_{\vec{\lambda}}(\mathcal{F}_Y) .
  \end{aligned}
\]
此外，(4.5) 与 (4.1) 中定义的 $\hat{\mathfrak{g}}_N(\mathcal{B})$ 与
$\hat{\mathfrak{g}}_N(Y)$ 在 $\cH_{\vec{\lambda}}(\mathcal{B})$、
$\cH_{\vec{\lambda}}^\dagger(\mathcal{B})$、$\cH_{\vec{\lambda}}(Y)$ 与
$\cH_{\vec{\lambda}}^\dagger(Y)$ 上的作用，以及 $\hat{\mathfrak{g}}_N$ 在
$\cH_{\vec{\lambda}}$ 上的作用，关于上述典范同构都是相容的。
\end{lemma}
% [Tutor Note: 引理 4.3 的说法是「纤维化与取纤维可交换」的代数版本：
%  $\C_s=\cO_{\mathcal{B},s}/\mathfrak{m}_s$ 把 $\mathcal{B}$ 上的层拉回到
%  点 $s$，而「先张量 $\C_s$」与「先取 $s$ 上的纤维」在四个层上给出同一个
%  对象；证明的关键是第三条（张量积与 $\pi_*$ 的交换），见下页的基变换
%  （base change），其余几条由定义直接得到。紧随其后的「更一般地」一段把
%  「点」换成任意全纯映射 $f:Y\to\mathcal{B}$，此时 $\C_s$ 换成 $\cO_Y$、
%  $\mathcal{F}_s$ 换成拉回族 $\mathcal{F}_Y$；$Y$ 可以不是流形（可以是带
%  幂零元素的方案），这正是后文讨论「形式形变」时需要的普遍性。原文 p.74–75]

\noindent\textit{证明.}\quad 第一、第二和第五个同构由定义是显然的。若 $k$ 充分大，
我们总有基变换（base change）
\[
  \begin{aligned}
  \C_s\otimes_{\cO_\mathcal{B}}\pi_*(\cO_\mathcal{C}(kS))
  &\simeq H^0\bigl(C_s,\cO_{C_s}\bigl(k\sum_{j=1}^Ns_j(s)\bigr)\bigr),\\
  \cO_Y\otimes_{\cO_\mathcal{B}}\pi_*(\cO_\mathcal{C}(kS))
  &\simeq H^0\bigl(C_Y,\cO_{C_Y}\bigl(k\sum_{j=1}^Ns_j(Y)\bigr)\bigr),
  \end{aligned}
\]
因为对每个点 $s\in\mathcal{B}$ 都有
\[
  H^1\bigl(C_s,\cO_{C_s}\bigl(k\sum_{j=1}^Ns_j(s)\bigr)\bigr)=0
\]
（例如参见 [Ha, Chap. III, Corollary 12.9] 或 [BS, Chap. III Corollary
3.5]）。这就蕴涵了第三个同构。
% [Tutor Note: 为什么需要 $H^1=0$：$\pi_*$ 与 $\otimes$ 的交换（基变换）
%  一般只给出谱序列，高阶项存活时就不是同构。这里
%  $\cO_{C_s}(k\sum_j s_j(s))$ 在亏格 $g$ 曲线上次数为 $kN$，$k$ 充分大时
%  次数 $>2g-2$，故 $H^1=0$（Serre 对偶 $H^1\cong H^0(K-k\sum s_j)^\vee$，
%  后者的次数 $2g-2-kN<0$）。原文把「$k$ 充分大」与引理 4.3 的陈述分开写，
%  实际上四个同构里的前两个与第三个都依赖这个 $k$ 的取法（通过
%  $\varinjlim_k$）。原文 p.75]

最后让我们考虑如下正合列（exact sequence）的交换图（commutative diagram）：
\[
  \begin{array}{ccc}
  & 0 & \\
  & \downarrow & \\
  \C_s\otimes\bigl(\hat{\mathfrak{g}}(\mathcal{F})
    \cH_{\vec{\lambda}}(\mathcal{B})\bigr)
  & \xrightarrow{\ \beta\ } &
  \bigl(\hat{\mathfrak{g}}(\mathcal{F})\otimes\C_s\bigr)
    \bigl(\C_s\otimes\cH_{\vec{\lambda}}(\mathcal{B})\bigr)\\
  \downarrow\alpha & & \downarrow\epsilon\\
  \C_s\otimes\cH_{\vec{\lambda}}(\mathcal{B})
  & \xrightarrow{\ \gamma\ } & \cH_{\vec{\lambda}}\\
  \downarrow & & \downarrow\\
  \C_s\otimes\cV_{\vec{\lambda}}(\mathcal{F})
  & \xrightarrow{\ \delta\ } & \cV_{\vec{\lambda}}(\mathcal{F}_s)\\
  \downarrow & & \downarrow\\
  0 & & 0
  \end{array}
\]
上述论证表明映射 $\beta$ 是满射（surjective）而映射 $\gamma$ 是同构。因此，
交换图的可交换性蕴涵映射 $\gamma$ 诱导出 $\operatorname{Im}(\alpha)$ 与
$\operatorname{Im}(\epsilon)$ 之间的同构。从而映射 $\delta$ 是同构。
\hfill$\square$
% [Tutor Note: 这段是引理 4.3 最后一条（$\cV$ 的纤维化）的证明骨架，值得补成
%  完整的追图：两列正合列分别是
%  $0\to\C_s\otimes(\hat{\mathfrak{g}}\cH)\to\C_s\otimes\cH
%   \to\C_s\otimes\cV\to0$
%  与 $0\to(\hat{\mathfrak{g}}\otimes\C_s)(\C_s\otimes\cH)\to\cH_{\vec\lambda}
%   \to\cV(\mathcal{F}_s)\to0$。$\beta$ 满、$\gamma$ 同构，于是
%  $\gamma(\operatorname{Im}\alpha)=\operatorname{Im}\epsilon$，两端的
%  $\operatorname{Im}$ 都被同构搬运，剩下的 $\delta$ 只能也是同构（蛇引理）。
%  注意上下两列的正合性用到了 $\cV$ 与 $\cH$ 的平坦性之外的额外信息：
%  这里 $\C_s$ 是域，张量积是右正合的，故上列在 $\C_s\otimes\cH$ 处正合
%  需要 $\hat{\mathfrak{g}}(\mathcal{F})\cH$ 的嵌入，原文默认了这一层。原文 p.75]

% [原文 p.76]
接下来我们将证明层 $\cV_{\vec{\lambda}}(\mathcal{F})$ 是凝聚
$\cO_\mathcal{B}$-模。首先我们引入一个滤过 $\{F_\bullet\}$，它是 (3.16) 与
(3.17) 中所给滤过的推广。$\cO_\mathcal{B}((\xi_j))$ 上的滤过
$\{F_\bullet\}$ 由
\[
  F_p\cO_\mathcal{B}((\xi_j)) = \cO_\mathcal{B}[[\xi_j]]\xi_j^{-p},
  \qquad p\in\Z
\]
定义。$\hat{\mathfrak{g}}_N(\mathcal{B})$ 上的滤过 $\{F_\bullet\}$ 由
\begin{equation}
  F_p\hat{\mathfrak{g}}_N(\mathcal{B}) =
  \begin{cases}
    \mathfrak{g}\otimes\bigoplus_{j=1}^NF_p\cO_\mathcal{B}((\xi_j))
    \oplus\cO_\mathcal{B}\cdot c, & p\ge 0,\\[2pt]
    \mathfrak{g}\otimes\bigoplus_{j=1}^NF_p\cO_\mathcal{B}((\xi_j)), & p<0.
  \end{cases}
  \tag{4.7}
\end{equation}
定义。$\cH_{\vec{\lambda}}(\mathcal{B})$ 上的滤过 $\{F_\bullet\}$ 由
\begin{equation}
  F_p\cH_{\vec{\lambda}}(\mathcal{B})
  = \cO_\mathcal{B}\otimes F_p\cH_{\vec{\lambda}} \tag{4.8}
\end{equation}
定义，其中 $F_p\cH_{\vec{\lambda}}
=F_p(U(\hat{\mathfrak{g}}_N))\cdot V_{\vec{\lambda}}$（见 (3.17)）。定义
\begin{equation}
  W_m(\mathcal{B}) = \cH_{\vec{\lambda}}(\mathcal{B})
  \big/\bigl(\mathfrak{g}[1]\oplus\hat{\mathfrak{g}}_N^{(m)}\bigr)
  \cH_{\vec{\lambda}}(\mathcal{B})
  = W_m\otimes\cO_\mathcal{B}, \tag{4.9}
\end{equation}
其中 $\mathfrak{g}[1]$ 与 $\hat{\mathfrak{g}}_N^{(m)}$ 分别由 (3.13) 与 (3.14)
定义，$W_m$ 由 (3.15) 定义。我们考虑
$\hat{\mathfrak{g}}_N(\mathcal{B})$ 的各个子代数上的诱导滤过，以及
$\cV_{\vec{\lambda}}(\mathcal{F})$ 与 $W_m(\mathcal{B})$ 上的商滤过，并把相应
的分次对象记为 $\operatorname{gr}^F_\bullet(\ )$。于是我们有如下分次
$\cO_\mathcal{B}$-模的正合列：
\[
  \begin{aligned}
  0&\longrightarrow\operatorname{gr}^F_\bullet
    (\hat{\mathfrak{g}}(\mathcal{F}))
  \longrightarrow\operatorname{gr}^F_\bullet\cH_{\vec{\lambda}}(\mathcal{B})
  \longrightarrow\operatorname{gr}^F_\bullet\cV_{\vec{\lambda}}(\mathcal{F})
  \longrightarrow 0,\\
  0&\longrightarrow\operatorname{gr}^F_\bullet\bigl(
    \bigl(\cO_\mathcal{B}\otimes_\C(\hat{\mathfrak{g}}_N^{(m)}
    \oplus\mathfrak{g}[1])\bigr)\cdot\cH_{\vec{\lambda}}(\mathcal{B})\bigr)
  \longrightarrow\operatorname{gr}^F_\bullet\cH_{\vec{\lambda}}(\mathcal{B})
  \longrightarrow\operatorname{gr}^F_\bullet W_m(\mathcal{B})
  \longrightarrow 0 .
  \end{aligned}
\]
由 Riemann-Roch 定理，对任意整数 $m\ge 2g-1$ 我们有
\[
  \operatorname{gr}^F_\bullet\bigl(\cO_\mathcal{B}\otimes_\C
  (\hat{\mathfrak{g}}_N^{(m)}\oplus\mathfrak{g}[1])\bigr)
  \subset\operatorname{gr}^F_\bullet(\hat{\mathfrak{g}}(\mathcal{F})) .
\]
固定这样的 $m$。于是由 (3.18) 与上述正合列，我们得到
$\cO_\mathcal{B}$-模的满同态
\[
  \operatorname{gr}^F_\bullet W_m(\mathcal{B})
  \to\operatorname{gr}^F_\bullet\cV_{\vec{\lambda}}(\mathcal{F})\to 0 .
\]
由命题 3.7，向量空间 $W_m$ 是有限维的。因此由 (4.9)，
$\operatorname{gr}^F_\bullet W_m(\mathcal{B})$ 是凝聚
$\cO_\mathcal{B}$-模。从而 $\cV_{\vec{\lambda}}(\mathcal{F})$ 也是凝聚
$\cO_\mathcal{B}$-模。这样我们就证明了如下定理。

\begin{theorem}[4.4]
层 $\cV_{\vec{\lambda}}(\mathcal{F})$ 与
$\cV_{\vec{\lambda}}^\dagger(\mathcal{F})$ 都是凝聚
$\cO_\mathcal{B}$-模。
\end{theorem}
% [Tutor Note: 这段证明的骨架值得单独理一遍：先用 (4.7)(4.8) 把「按 $\xi_j$
%  的极点阶数」的滤过搬到层上，再注意 $\hat{\mathfrak{g}}(\mathcal{F})$ 的
%  分次与 $\cO_\mathcal{B}\otimes(\hat{\mathfrak{g}}_N^{(m)}\oplus\mathfrak{g}[1])$
%  的分次在 $m\ge 2g-1$ 时相容（Riemann-Roch 给出维数比较），于是
%  $\operatorname{gr}^F_\bullet\cV$ 成为有限维空间 $W_m$ 的分次商层的商，
%  有限维向量空间张量 $\cO_\mathcal{B}$ 显然是凝聚的，商的核与像都是凝聚的，
%  故 $\operatorname{gr}^F_\bullet\cV$ 凝聚；最后用「分次凝聚 ⟹ 原层凝聚」
%  （滤过有界、$\cO_\mathcal{B}$ 诺特）回到 $\cV$ 本身。原文一句
%  "Hence $\cV_{\vec{\lambda}}(\mathcal{F})$ is also a coherent..." 跳过了
%  最后这一步，是本块最容易漏的一环。原文 p.76]
% [Tutor Note: 定理 4.4 同时断言 $\cV^\dagger$ 凝聚：这一步用对偶
%  $\cV^\dagger=\underline{\operatorname{Hom}}_{\cO_\mathcal{B}}
%  (\cV,\cO_\mathcal{B})$ 与「凝聚层之间的 $\underline{\operatorname{Hom}}$
%  仍凝聚」；但需 $\cV$ 凝聚且 $\cO_\mathcal{B}$ 诺特，$\mathcal{B}$ 是
%  有限维复流形（故其结构层诺特）。原文未展开。原文 p.76]

\subsection{局部自由性 I（光滑情形）}

本节我们将证明共形块层
$\cV_{\vec{\lambda}}^\dagger(\mathcal{F})$ 的局部自由性（local freeness）。
为此，首先我们引入 $\cO_\mathcal{B}$-子模
$\mathcal{L}(\mathcal{F})\subset\bigoplus_{j=1}^N
\cO_\mathcal{B}((\xi_j^{-1}))\dfrac{d}{d\xi_j}$ 以及
$\mathcal{L}(\mathcal{F})$ 作为一阶扭转微分算子（first order twisted
differential operator）在真空层与余真空层上的作用。在 §5.1 中，这一作用将
被用来在真空层上定义射影平坦联络（projectively flat connection）。
% [Tutor Note: $\cO_\mathcal{B}((\xi_j^{-1}))$ 是「以 $\xi_j^{-1}$ 为变量的
%  洛朗级数」，即允许 $\xi_j$ 的任意阶极点、只允许有限个正幂的级数；乘上
%  $\frac{d}{d\xi_j}$ 后，这些正是「在标点 $s_j$ 的形式邻域里允许极点的
%  向量场」。$\mathcal{L}(\mathcal{F})$ 是它的一个 $\cO_\mathcal{B}$-子模，
%  由那些「能整体定义在族上」的向量场组成——下一节会看到它由
%  $\hat{\mathfrak{g}}(\mathcal{F})$ 的 Sugawara 型元素给出，因而其维数
%  与 $W_m$ 一样是有限的。原文 p.76]

% [原文 p.77]
\noindent\textbf{4.2.1 $\cV_{\vec{\lambda}}(\mathcal{F})$ 上的微分算子.}\quad
设 $\mathcal{F}^{(0)}=(\pi:\mathcal{C}\to\mathcal{B};s_1,\dots,s_N)$ 是亏格为
$g$ 的 $N$-点稳定曲线的\emph{通用族}（versal family）。以 $\Sigma$ 记
$\mathcal{F}^{(0)}$ 的纤维的二重点（double point）的轨迹，并置
$D=\pi(\Sigma)$。注意 $\Sigma$ 是 $\mathcal{C}$ 中余维（codimension）为二的
非奇异子流形（submanifold），而 $D$ 是 $\mathcal{B}$ 中的除子，其不可约分支
（irreducible component）$D_i$，$i=1,2,\dots,k$，都是非奇异的（见引理
1.29）。假定给定形式坐标
\[
  \eta_j : \hat{\cO}_{\mathcal{C}/s_j(\mathcal{B})}\simeq
  \cO_\mathcal{B}[[\xi]],\qquad j=1,2,\dots,N .
\]
为简单起见，下文我们假定 $\eta_j^{-1}(\xi)$ 在 $s_j(\mathcal{B})$ 的一个邻域
内全纯。一般情形可以用全纯坐标去逼近形式坐标来处理。我们自由地使用如下记号：
\[
  S_j=s_j(\mathcal{B}),\qquad S=\sum_{j=1}^NS_j,\qquad
  \xi_j=\eta_j^{-1}(\xi) .
\]
于是 $\mathcal{F}=(\pi:\mathcal{C}\to\mathcal{B};s_1,\dots,s_N;
\eta_1,\dots,\eta_N)$ 是带形式邻域的亏格 $g$ 的 $N$-点稳定曲线族。对每个
$\vec{\lambda}\in(P_\ell)^{\otimes N}$，我们可以定义真空层
$\cV_{\vec{\lambda}}^\dagger(\mathcal{F})$ 与余真空层
$\cV_{\vec{\lambda}}(\mathcal{F})$。它们分别是
$\cH_{\vec{\lambda}}^\dagger(\mathcal{F})$ 的子层与
$\cH_{\vec{\lambda}}(\mathcal{F})$ 的商层。

首先回忆，我们有如下 $\cO_\mathcal{B}$-模的正合列（见 (1.53)）：
\begin{equation}
  0\to\pi_*(\Theta_{\mathcal{C}/\mathcal{B}}(*S))
  \xrightarrow{\ b\ }\bigoplus_{j=1}^N\cO_\mathcal{B}[\xi_j^{-1}]
  \frac{d}{d\xi_j}
  \xrightarrow{\ \theta\ }R^1\pi_*\Theta_{\mathcal{C}/\mathcal{B}}(-S)
  \to 0 \tag{4.10}
\end{equation}
更确切地说，对任意正整数 $m\ge 4g-3$，我们有如下正合列：
\begin{equation}
  0\to\pi_*(\Theta_{\mathcal{C}/\mathcal{B}}(mS))
  \xrightarrow{\ b_m\ }\bigoplus_{j=1}^N\bigoplus_{k=0}^m
  \cO_\mathcal{B}\xi_j^{-k}\frac{d}{d\xi_j}
  \xrightarrow{\ \theta_m\ }R^1\pi_*\Theta_{\mathcal{C}/\mathcal{B}}(-S)
  \to 0 \tag{4.11}
\end{equation}
注意映射 $b$ 与 $b_m$ 对应着关于 $\xi_j$ 的、截断到零阶的洛朗展开。为了定义
作用在真空层与余真空层上的一阶扭转微分算子，我们需要按如下方式修改正合列
(4.10)。
存在正合列
\[
  0\to\Theta_{\mathcal{C}/\mathcal{B}}\to\Theta_\mathcal{C}
  \xrightarrow{\ d\pi\ }\pi^*\Theta_\mathcal{B}\to 0 ,
\]
其中 $\Theta_{\mathcal{C}/\mathcal{B}}$ 是与 $\pi$ 的纤维相切的向量场
（vector field）层。置
\[
  \Theta'_{\mathcal{C},\pi}
  = d\pi^{-1}\bigl(\pi^{-1}\Theta_\mathcal{B}(-\log D)\bigr) .
\]
于是 $\Theta'_{\mathcal{C},\pi}$ 是 $\mathcal{C}$ 上沿 $\Sigma$ 相切、且其
水平分量（horizontal component）在 $\pi$ 的纤维上为常数的向量场层。也就是说，
$\Theta'_{\mathcal{C},\pi}$ 由如下形状的全纯向量场的芽（germ）组成：
\[
  a(z,u)\frac{\partial}{\partial z}
  +\sum_{i=1}^mb_i(u)u_i\frac{\partial}{\partial u_i}
  +\sum_{i=m+1}^nb_i(u)\frac{\partial}{\partial u_i},
\]
其中 $(z,u_1,\dots,u_n)$ 是局部坐标（local coordinate）组，使得映射 $\pi$
表示成投影（projection）
\[
  \pi(z,u_1,\dots,u_n)=(u_1,\dots,u_n) .
\]
% [Tutor Note: (4.10)(4.11) 是本节的技术核心：它把「$\mathcal{C}$ 上在 $S$ 处
%  允许极点的向量场」实现为「每个标点处 $\xi_j$ 的洛朗级数乘 $\frac{d}{d\xi_j}$」。
%  $b$ 就是逐点作洛朗展开（与 §1 的 $\tilde t$ 同源），$\theta$ 是「取留数」
%  型的障碍映射，落在 $R^1\pi_*$ 中；右端正合给出「单值化障碍为零」的判据。
%  (4.11) 是它的有限阶截断版，$m\ge4g-3$ 保证 $\theta_m$ 仍满（其余项被
%  高阶极点吸收），这是下文把 $\mathcal{L}(\mathcal{F})$ 取成有限秩层的依据。
%  原文 p.77]
% [Tutor Note: $\Theta'_{\mathcal{C},\pi}$ 的定义 $d\pi^{-1}(\pi^{-1}
%  \Theta_\mathcal{B}(-\log D))$ 意思是：把 $\mathcal{B}$ 上「沿 $D$ 允许
%  对数极点」的向量场沿 $\pi$ 拉回，再在 $\mathcal{C}$ 上取原像。条件
%  「水平分量沿纤维为常数」正是对数极点的坐标表现：$u_i$ 方向的系数对
%  $i\le m$ 带因子 $u_i$（$u_1=\cdots=u_m=0$ 是二重点 $\Sigma$ 的局部方程，
%  故 $u_i\partial_{u_i}$ 在 $\Sigma$ 上不消失、且给出对数极点），对
%  $i>m$ 则不带因子（那些方向是 $\Sigma$ 上的光滑方向）。原文 p.77]

% [原文 p.78]
而 $\pi(\Sigma)=D$ 由方程
\[
  u_1\cdot u_2\cdots u_m=0
\]
给出。

更一般地，我们可以定义层 $\Theta'_\mathcal{C}(mS)_\pi$ 为如下形状的亚纯向量场
的芽所组成的层：
\[
  A(z,u)\frac{\partial}{\partial z}
  +\sum_{i=1}^mB_i(u)u_i\frac{\partial}{\partial u_i}
  +\sum_{i=m+1}^nB_i(u)\frac{\partial}{\partial u_i},
\]
其中 $A(z,u)$ 沿 $S$ 的极点阶数至多为 $m$。现在我们有正合列
\begin{equation}
  0\to\Theta_{\mathcal{C}/\mathcal{B}}(mS)
  \to\Theta'_\mathcal{C}(mS)_\pi\xrightarrow{\ d\pi\ }
  \pi^{-1}\Theta_\mathcal{B}(-\log D)\to 0 . \tag{4.12}
\end{equation}
注意 $\Theta'_\mathcal{C}(mS)_\pi$ 借助向量场的通常括号运算具有李代数层的
结构，而上述正合列正是李代数层的正合列。

对 $m>\dfrac{1}{N}(2g-2)$，由 (4.12) 我们得到 $\cO_\mathcal{B}$-模的正合列
\begin{equation}
  0\to\pi_*\Theta_{\mathcal{C}/\mathcal{B}}(mS)
  \to\pi_*\Theta'_\mathcal{C}(mS)_\pi\xrightarrow{\ d\pi\ }
  \Theta_\mathcal{B}(-\log D)\to 0 , \tag{4.13}
\end{equation}
它同时也是李代数层的正合列。令 $m\to\infty$，我们得到正合列
\begin{equation}
  0\to\pi_*\Theta_{\mathcal{C}/\mathcal{B}}(*S)
  \to\pi_*\Theta'_\mathcal{C}(*S)_\pi\xrightarrow{\ d\pi\ }
  \Theta_\mathcal{B}(-\log D)\to 0 . \tag{4.14}
\end{equation}

正合列 (4.10) 与 (4.12) 由如下交换图相联系：
\[
  \begin{array}{ccccccc}
  0 & \to & \pi_*\Theta_{\mathcal{C}/\mathcal{B}}(*S) & \to &
  \pi_*\Theta'_\mathcal{C}(*S)_\pi & \xrightarrow{\ d\pi\ } &
  \Theta_\mathcal{B}(-\log D)\to 0\\[2pt]
  & & \| & & \downarrow p & & \downarrow\rho\\[2pt]
  0 & \to & \pi_*\Theta_{\mathcal{C}/\mathcal{B}}(*S) & \to &
  \bigoplus_{j=1}^N\cO_\mathcal{B}[\xi_j^{-1}]\dfrac{d}{d\xi_j}
  & \xrightarrow{\ \theta\ } &
  R^1\pi_*\Theta_{\mathcal{C}/\mathcal{B}}(-S)\to 0
  \end{array}
\]
其中 $\rho$ 是族 $\mathcal{F}^{(0)}$ 的小平–Spencer 映射
（Kodaira-Spencer mapping），而 $p$ 由「取
$\pi_*\Theta_\mathcal{C}(mS)_\pi$ 中的向量场在 $s_j(\mathcal{B})$ 处的洛朗
展开的 $\dfrac{d}{d\xi_j}$ 部分的非正部分」给出。因为我们的族
$\mathcal{F}^{(0)}$ 是通用的，小平–Spencer 映射 $\rho$ 是
$\cO_\mathcal{B}$-模的同构。因此 $p$ 是同构。令
\[
  \tilde{p} : \pi_*\Theta'_\mathcal{C}(*S)_\pi\to\bigoplus_{j=1}^N
  \cO_\mathcal{B}((\xi_j))\frac{d}{d\xi_j}
\]
为「取在 $s_j(\mathcal{B})$ 处的洛朗展开的 $\dfrac{d}{d\xi_j}$ 部分」所得的
映射。由于 $p$ 是同构，$\tilde{p}$ 是单射。置
\[
  \mathcal{L}(\mathcal{F}):=\tilde{p}\bigl(\pi_*\Theta'_\mathcal{C}(*S)_\pi
  \bigr) .
\]
于是我们有正合列
\begin{equation}
  0\to\pi_*\Theta_{\mathcal{C}/\mathcal{B}}(*S)\to\mathcal{L}(\mathcal{F})
  \xrightarrow{\ \theta\ }\Theta_\mathcal{B}(-\log D)\to 0 . \tag{4.15}
\end{equation}
% [Tutor Note: 这一段是「局部自由性 I」的全部技术装置，值得把逻辑链写清：
%  (i) 上半行 (4.14) 是「$\mathcal{C}$ 上沿 $\Sigma$ 对数、在 $S$ 处允许极点的
%  向量场」到「$\mathcal{B}$ 上沿 $D$ 对数的向量场」的满射，核是沿纤维的向量场；
%  (ii) 下半行 (4.10) 把同一批向量场在标点处展开成洛朗级数；
%  (iii) $p$ 是「取非正部分」，$\rho$ 是 Kodaira-Spencer，二者构成方块；
%  (iv) 通用性使 $\rho$ 是同构（这正是「通用族」的定义：形变方向与
%  $H^1(\Theta)$ 同构），于是 $p$ 同构（方块追图）；
%  (v) $\tilde p$ 比 $p$ 多取了 $\xi_j$ 的正幂部分，故只是单射；
%  (vi) 把 $\tilde p$ 的像定义为 $\mathcal{L}(\mathcal{F})$，就得到 (4.15)。
%  关键在于 (v)→(vi)：$\mathcal{L}(\mathcal{F})$ 从而是有限秩局部自由层，
%  这是把它当作「联络」的系数层的理由。原文 p.78]
% [Tutor Note: 原文在此处印 $\pi_*\Theta_\mathcal{C}(mS)_\pi$，按上下文（它应当
%  是上文的 $\pi_*\Theta'_\mathcal{C}(mS)_\pi$，因为只有 $\Theta'$ 才含 $\Sigma$
%  的对数极点、才在 $d\pi$ 下满射到 $\Theta_\mathcal{B}(-\log D)$）应为
%  $\pi_*\Theta'_\mathcal{C}(mS)_\pi$。照原文排，见回报 ERRATA。原文 p.78]

% [原文 p.79]
$\cO_\mathcal{B}$-模的正合列。$\mathcal{L}(\mathcal{F})$ 的李括号
$[\ ,\ ]_d$ 由 $\pi_*\Theta'_\mathcal{C}(*S)_\pi$ 的李括号经映射 $\tilde{p}$
得到。于是对 $\vec{\ell},\vec{m}\in\mathcal{L}(\mathcal{F})$ 我们有
\begin{equation}
  [\vec{\ell},\vec{m}]_d
  = [\vec{\ell},\vec{m}]_0+\theta(\vec{\ell})(\vec{m})
    -\theta(\vec{m})(\vec{\ell}) \tag{4.16}
\end{equation}
这里 $[\ ,\ ]_0$ 是形式向量场的通常括号，而 $\theta(\vec{\ell})$ 对
\[
  \vec{m}=\Bigl(m_1\frac{d}{d\xi_1},\dots,m_N\frac{d}{d\xi_N}\Bigr)
\]
的作用由
\[
  \Bigl(\theta(\vec{\ell})(m_1)\frac{d}{d\xi_1},\dots,
  \theta(\vec{\ell})(m_N)\frac{d}{d\xi_N}\Bigr)
\]
定义。此时正合列 (4.15) 也是李代数层的正合列。
% [Tutor Note: (4.16) 的修正项 $\theta(\vec\ell)(\vec m)-\theta(\vec m)(\vec\ell)$
%  来自「$\mathcal{L}(\mathcal{F})$ 的元素是 $\pi_*\Theta'_\mathcal{C}(*S)_\pi$
%  的元素沿 $\tilde p$ 的像，而 $\tilde p$ 丢掉了 $\xi_j$ 的正幂部分」：两个向量场
%  在 $\mathcal{C}$ 上先作括号、再展开，与先展开、再按形式向量场作括号，
%  相差的正是 $\theta$ 作用在系数上的项。所以 $\mathcal{L}(\mathcal{F})$ 的
%  李代数结构是 $\pi_*\Theta'_\mathcal{C}(*S)_\pi$ 的结构经 $\tilde p$ 的
%  「transport」；称 $[\ ,\ ]_d$ 的下标 $d$ 为「微分算子化」的括号。原文 p.79]

下面定义 $\mathcal{L}(\mathcal{F})$ 在 $\cH_{\vec{\lambda}}(\mathcal{B})$ 上
的作用。

\begin{definition}[4.5]
对 $\vec{\ell}=(\ell_1,\dots,\ell_N)\in\mathcal{L}(\mathcal{F})$，
$\cH_{\vec{\lambda}}(\mathcal{B})$ 上的作用 $D(\vec{\ell})$ 由
\begin{equation}
  D(\vec{\ell})(F\otimes|\Phi\rangle)
  = \theta(\vec{\ell})(F)\otimes|\Phi\rangle
  -F\cdot\Bigl(\sum_{j=1}^N\rho_j(T[\ell_j])|\Phi\rangle\Bigr)
  \tag{4.17}
\end{equation}
定义，其中
\[
  F\in\cO_\mathcal{B},\qquad |\Phi\rangle\in\cH_{\vec{\lambda}} .
\]
\end{definition}

下述命题是定义与引理 2.13 的容易推论。

\begin{proposition}[4.6]
如上定义的 $\vec{\ell}\in\mathcal{L}(\mathcal{F})$ 在
$\cH_{\vec{\lambda}}(\mathcal{B})$ 上的作用 $D(\vec{\ell})$ 具有如下性质。
\begin{enumerate}
\item[(1)] 对任意 $f\in\cO_\mathcal{B}$ 我们有
\[
  D(f\vec{\ell})=fD(\vec{\ell}) .
\]
\item[(2)] 对 $\vec{\ell},\vec{m}\in\mathcal{L}(\mathcal{F})$ 我们有
\[
  [D(\vec{\ell}),D(\vec{m})] = D([\vec{\ell},\vec{m}]_d)
  +\frac{c_v}{12}\sum_{j=1}^N\operatorname{Res}_{\xi_j=0}
  \Bigl(\frac{d^3\ell_j}{d\xi_j^3}m_jd\xi_j\Bigr)\cdot\mathrm{id} .
\]
\item[(3)] 对 $f\in\cO_\mathcal{B}$ 与
$|\phi\rangle\in\cH_{\vec{\lambda}}(\mathcal{B})$ 我们有
\[
  D(\vec{\ell})(f|\phi\rangle)
  =(\theta(\vec{\ell})(f))|\phi\rangle+fD(\vec{\ell})|\phi\rangle .
\]
\end{enumerate}
也就是说，若 $\theta(\vec{\ell})\ne 0$，则 $D(\vec{\ell})$ 是一阶扭转微分
算子。
\end{proposition}
% [Tutor Note: 这三条性质正是「$D$ 是射影平坦联络的候选」的代数检验：
%  (1) 说明 $D$ 是 $\cO_\mathcal{B}$-线性的（$\theta$ 是 $\cO_\mathcal{B}$-线性
%  的向量场作用，$T[\ell_j]$ 也随 $\ell_j$ 线性）；(2) 说明 $D$ 的曲率是
%  中心元——修正项 $\frac{c_v}{12}\sum_j\operatorname{Res}((d^3\ell_j/d\xi_j^3)
%  m_jd\xi_j)\cdot\mathrm{id}$ 恰是中心荷 $c_v$ 的倍数，故 $D$ 给出的是
%  \emph{射影}平坦联络而非平坦联络（这正是本节的结论）；(3) 说明 $D$ 对
%  函数是「一阶」的，即 $D$ 确实是一阶微分算子，$\theta(\vec\ell)$ 是它的
%  「符号」（principal symbol）。至于 (4.17) 右端第二项里的 $T[\ell_j]$：
%  它是 §2 的 Sugawara 型元素，把向量场 $\ell_j\frac{d}{d\xi_j}$ 通过
%  $\tilde t$ 的逆映到 $\cH$ 上的作用（$T[\ ]$ 是能量-动量张量的「抹平」）。
%  因此 $D$ 的两项分别对应「几何上沿 $\vec\ell$ 的方向导数」与「量子化后的
%  $T$ 插入」，一阶扭转微分算子的「扭转」就体现在第二项。原文 p.79]

我们用
\begin{equation}
  D(\vec{\ell})(F\otimes\langle\Psi|)
  = (\theta(\vec{\ell})F)\otimes\langle\Psi|
  +\sum_{j=1}^NF\cdot\bigl(\langle\Psi|\rho_j(T[\ell_j])\bigr)
  \tag{4.18}
\end{equation}
定义 $\mathcal{L}(\mathcal{F})$ 在 $\cH_{\vec{\lambda}}^\dagger(\mathcal{B})$
上的对偶作用，其中
\[
  F\in\cO_\mathcal{B},\qquad
  \langle\Psi|\in\cH_{\vec{\lambda}}^\dagger(\mathcal{B}) .
\]
于是对任意 $|\tilde{\Phi}\rangle\in\cH_{\vec{\lambda}}(\mathcal{B})$ 与
$\langle\tilde{\Psi}|\in\cH_{\vec{\lambda}}^\dagger(\mathcal{B})$，我们有
\begin{equation}
  \{D(\vec{\ell})\langle\tilde{\Psi}|\}|\tilde{\Phi}\rangle
  +\langle\tilde{\Psi}|\{D(\vec{\ell})|\tilde{\Phi}\rangle\}
  = \theta(\vec{\ell})\langle\tilde{\Psi}|\tilde{\Phi}\rangle . \tag{4.19}
\end{equation}
这与对偶联络的通常定义一致。另见下文命题 4.12 与 (4.23)。

% [原文 p.80]
下面我们要说明算子 $D(\vec{\ell})$ 作用在
$\cV_{\vec{\lambda}}(\mathcal{F})$ 上。

\begin{proposition}[4.7]
对任意 $\vec{\ell}\in\mathcal{L}(\mathcal{F})$ 我们有
\[
  D(\vec{\ell})\bigl(\hat{\mathfrak{g}}(\mathcal{F})
  \cH_{\vec{\lambda}}(\mathcal{B})\bigr)
  \subset\hat{\mathfrak{g}}(\mathcal{F})
  \cH_{\vec{\lambda}}(\mathcal{B}) .
\]
因此 $D(\vec{\ell})$ 作用在 $\cV_{\vec{\lambda}}(\mathcal{F})$ 上。此外，若
$\theta(\vec{\ell})\ne 0$，它是一阶微分算子。
\end{proposition}

\noindent\textit{证明.}\quad
$\hat{\mathfrak{g}}(\mathcal{F})\cH_{\vec{\lambda}}(\mathcal{B})$ 的元素是
如下形状的元素的线性组合：
\[
  F\otimes\Bigl(\sum_{j=1}^N\rho_j(X\otimes t_j(h))|\Phi\rangle\Bigr),
\]
其中
\[
  F\in\cO_\mathcal{B},\quad X\in\mathfrak{g},\quad
  h\in\pi_*\cO_\mathcal{C}(*S),\quad
  |\Phi\rangle\in\cH_{\vec{\lambda}},
\]
而 $t_j(h)$ 是 $h$ 在 $S_j=s_j(\mathcal{B})$ 处关于参数 $\xi_j$ 的洛朗展开。
首先我们要把下述等式作为 $\cH_{\vec{\lambda}}(\mathcal{B})$ 上的算子来证明：
\begin{equation}
  \Bigl[D(\vec{\ell}),\sum_{j=1}^N\rho_j(X\otimes t_j(h))\Bigr]
  = \sum_{j=1}^N\rho_j\Bigl(X\otimes
    \bigl\{\theta(\vec{\ell})(t_j(h))+\ell_j(t_j(h))\bigr\}\Bigr),
  \tag{4.20}
\end{equation}
其中 $\theta(\vec{\ell})$ 作用在 $t_j(h)$ 的系数上。由命题 4.6 的 (3)，只需把
等式 (4.20) 作为 $\cH_{\vec{\lambda}}$ 上的算子来证明。对
$|\Phi\rangle\in\cH_{\vec{\lambda}}$，由 (2.19) 我们有
\[
  \begin{aligned}
  &D(\vec{\ell})\Bigl(\sum_{j=1}^N\rho_j(X\otimes t_j(h))|\Phi\rangle\Bigr)
  -\sum_{j=1}^N\rho_j(X\otimes t_j(h))\bigl(D(\vec{\ell})|\Phi\rangle\bigr)\\
  &\quad = \sum_{j=1}^N\bigl\{\rho_j(X\otimes\theta(\vec{\ell})(t_j(h)))
    -T[\ell_j]\rho_j(X\otimes t_j(h))\bigr\}|\Phi\rangle\\
  &\qquad +\sum_{j=1}^N\rho_j(X\otimes t_j(h))T[\ell_j]|\Phi\rangle\\
  &\quad = \sum_{j=1}^N\bigl\{\rho_j(X\otimes\theta(\vec{\ell})(t_j(h)))
    +\rho_j(X\otimes\ell_j(t_j(h)))\bigr\}|\Phi\rangle .
  \end{aligned}
\]
但 $\theta(\vec{\ell})(t_j(h))+\ell_j(t_j(h))$ 不是别的，正是某个亚纯函数
（meromorphic function）$\tau(h)$ 在 $s_j(\mathcal{B})$ 处的洛朗展开，其中
$\tau=\tilde{p}^{-1}(\vec{\ell})\in\pi_*(\Theta'_\mathcal{C}(*S)_\pi)$。
于是我们得到所要的结果。\hfill$\square$
% [Tutor Note: 这段计算的关键是 (2.19)：$[T[\ell_j],\rho_j(X\otimes t_j(h))]
%  =\rho_j(X\otimes\ell_j(t_j(h)))$。把它展开成
%  $T[\ell_j]\rho_j(X\otimes t_j)=\rho_j(X\otimes t_j)T[\ell_j]
%  +\rho_j(X\otimes\ell_j(t_j))$ 之后，与 (4.17) 里 $D(\vec\ell)$ 的
%  $-\sum_j\rho_j(T[\ell_j])$ 项相遇，$\pm T[\ell_j]$ 恰好消去，剩下的正是
%  $\theta(\vec\ell)(t_j)+\ell_j(t_j)$。这也解释了 (4.20) 的右端为什么是
%  这两项之和：$\theta(\vec\ell)$ 来自几何方向导数经展开后的常数项修正，
%  $\ell_j$ 来自 $T$ 与 $\rho_j$ 的对易子。原文 p.80]

同样可以说明 $D(\vec{\ell})$ 作用在
$\cV_{\vec{\lambda}}^\dagger(\mathcal{F})$ 上。这一事实的另一证明可在下文
命题 4.12 的证明中找到。

\noindent\textbf{4.2.2 光滑族的局部自由性.}\quad

\begin{theorem}[4.8]
$\cO_\mathcal{B}$-模 $\cV_{\vec{\lambda}}(\mathcal{F})$ 在
$\mathcal{B}\setminus D$ 上是局部自由的（locally free），其中 $D$ 是由
“纤维为奇曲线（singular curve）的点”组成的底空间 $\mathcal{B}$ 的轨迹。
\end{theorem}

% [原文 p.81]
\noindent\textit{证明.}\quad 由引理 4.3，对任意点
$x\in\mathcal{B}\setminus D$ 我们有同构
\[
  \C_x\otimes_{\cO_\mathcal{B}}\cV_{\vec{\lambda}}(\mathcal{F})
  \simeq\cV_{\vec{\lambda}}(\mathcal{F}_x),
\]
其中我们置 $\C_x=\cO_{\mathcal{B}}/\mathfrak{m}_x$，且
\[
  \mathcal{F}_x := (\mathcal{C}_x=\pi^{-1}(x),s_1(x),\dots,s_N(x);
  \eta_1|_{\mathcal{C}_x},\dots,\eta_N|_{\mathcal{C}_x}) .
\]
设 $v_1,\dots,v_m$ 是 $\cV_{\vec{\lambda}}(\mathcal{F})$ 在 $x$ 的一个邻域内
的局部全纯截面（local holomorphic section），使得
$\{v_1(x),\dots,v_m(x)\}$ 构成 $\cV_{\vec{\lambda}}(\mathcal{F}_x)$ 的基。
假定存在非平凡关系
\begin{equation}
  a_1v_1+a_2v_2+\cdots+a_mv_m=0, \tag{4.21}
\end{equation}
其中 $a_i$ 是 $x$ 的某邻域内的全纯函数。按我们的假定，
\[
  a_1(x)=a_2(x)=\cdots=a_m(x)=0 .
\]
必要时改变下标顺序，我们可以假定存在正整数 $k$ 使得
\[
  \begin{aligned}
  a_1 &\in \mathfrak{m}_x^k\setminus\mathfrak{m}_x^{k+1},\\
  a_i &\in \mathfrak{m}_x^l,\quad l\ge k,\quad i=2,3,\dots,m .
  \end{aligned}
\]
我们这样选取 $a_i$，使得在关系 (4.21) 中正整数 $k$ 是最小的。设 $\tau$ 是
$x$ 的某邻域内处处非零的局部全纯向量场，使得
$\tau(a_1)\in\mathfrak{m}_x^{k-1}$。存在 $\vec{\ell}\in\mathcal{L}(\mathcal{F})$
使 $\theta(\vec{\ell})=\tau$。把 $\tau=\theta(\vec{\ell})$ 作用于等式
(4.21)，我们得到等式
\begin{equation}
  \sum_{i=1}^m\Bigl(\tau(a_i)+\sum_{j=1}^ma_j\alpha_{ji}\Bigr)v_i=0,
  \tag{4.22}
\end{equation}
其中
\[
  D(\vec{\ell})(v_i)=\sum_{j=1}^m\alpha_{ij}v_j .
\]
于是关系 (4.22) 是非平凡的，且
\[
  \tau(a_1)+\sum_{j=1}^m\alpha_{j1}\in\mathfrak{m}_x^{k-1} .
\]
这与我们的假定矛盾。因此 $v_1,\dots,v_m$ 在 $\cO_\mathcal{B}$ 上线性无关。
从而 $\cV_{\vec{\lambda}}(\mathcal{F})$ 在 $x$ 处是局部自由的
$\cO_\mathcal{B}$-模。\hfill$\square$
% [Tutor Note: 证明的反证逻辑要抓住三点：(i) 若在 $x$ 处有关系，取 $k$ 为
%  所有非平凡关系中「最小阶」；(ii) $\tau$ 的取法使 $\tau(a_1)$ 正好降一阶
%  （因为 $\tau$ 非零且 $a_1\in\mathfrak{m}^k\setminus\mathfrak{m}^{k+1}$，
%  取 $\tau$ 沿 $a_1$ 的首项方向即可，这就是 $\mathfrak{m}^{k-1}$ 的来历）；
%  (iii) 作用后新关系 (4.22) 的首项系数比原来的 $k$ 更低一阶，故它非零，
%  却又给出一个关系，矛盾于 $k$ 的最小性。中间用到 $D(\vec\ell)(v_i)$ 是
%  $v_1,\dots,v_m$ 的 $\cO_\mathcal{B}$-线性组合（即 $D$ 保持局部自由层），
%  这一步原文省略，由命题 4.7 保证。另外「存在 $\vec\ell$ 使
%  $\theta(\vec\ell)=\tau$」用到正合列 (4.15)：
%  $0\to\pi_*\Theta_{\mathcal{C}/\mathcal{B}}(*S)\to
%  \mathcal{L}(\mathcal{F})\xrightarrow{\theta}\Theta_\mathcal{B}(-\log D)\to0$，
%  而 $\tau$ 在 $\mathcal{B}\setminus D$ 上全纯，故属于
%  $\Theta_\mathcal{B}(-\log D)$ 的截面，可沿 $\theta$ 提升；原文只说
%  "There exists $\vec\ell$ with $\theta(\vec\ell)=\tau$"，跳过了这一步。
%  原文 p.81]

对一个凝聚 $\cO_\mathcal{B}$-模 $\mathcal{G}$，使其非局部自由的点所成的轨迹
$M$ 是 $\mathcal{B}$ 的余维至少为二的闭解析子集（closed analytic subset）。
因此我们有如下推论。

\begin{corollary}[4.9]
设 $W$ 是 $\mathcal{B}$ 上使 $\cV_{\vec{\lambda}}(\mathcal{F})$ 非局部自由的
极大子集。则 $W$ 是 $\mathcal{B}$ 的解析子集，且
\[
  W\subsetneqq D .
\]
\end{corollary}

由于我们定义
\[
  \cV_{\vec{\lambda}}^\dagger(\mathcal{F})
  =\underline{\operatorname{Hom}}_{\cO_\mathcal{B}}
  (\cV_{\vec{\lambda}}(\mathcal{F}),\cO_\mathcal{B}),
\]
我们有如下推论。

% [原文 p.82]
\begin{corollary}[4.10]
$\cV_{\vec{\lambda}}^\dagger(\mathcal{F})|_{\mathcal{B}\setminus D}$ 是局部
自由的 $\cO_\mathcal{B}$-模，且对 $\mathcal{B}\setminus D$ 的任意子簇
（subvariety）$Y$，我们有 $\cO_Y$-模同构
\[
  \cO_Y\otimes_{\cO_\mathcal{B}}\cV_{\vec{\lambda}}^\dagger(\mathcal{F})
  \simeq\cV_{\vec{\lambda}}^\dagger(\mathcal{F}|_Y) .
\]
\end{corollary}

这两个推论在证明一般情形的局部自由性时起关键作用。上述推论蕴涵如下定理。

\begin{theorem}[4.11]
若 $\mathcal{F}$ 是带形式邻域的 $N$-点光滑曲线族，使得由它诱导的
$N$-点光滑曲线族 $\mathcal{F}^{(0)}$ 不必是通用的，则
$\cV_{\vec{\lambda}}(\mathcal{F})$ 与
$\cV_{\vec{\lambda}}^\dagger(\mathcal{F})$ 都是局部自由
$\cO_\mathcal{B}$-模，并且互为对偶。
\end{theorem}
% [Tutor Note: 定理 4.11 与定理 4.8 的区别在于「族可以不是通用的」：定理 4.8
%  证明的是通用族在 $\mathcal{B}\setminus D$ 上的局部自由性，而一般族可以
%  通过「拉回自通用族」得到（这正是推论 4.10 的用处：沿子簇 $Y$ 拉回保持
%  局部自由），故定理 4.11 是把局部自由性从通用族搬到任意光滑族。
%  「互为对偶」指 $\cV^\dagger=\underline{\operatorname{Hom}}_{\cO_\mathcal{B}}
%  (\cV,\cO_\mathcal{B})$ 且自然的求值映射是同构——这是局部自由层的标准性质。
%  原文 p.82]

上述讨论的另一个重要推论如下。

\begin{proposition}[4.12]
对每个元素 $\vec{\ell}\in\mathcal{L}(\mathcal{F})$，$D(\vec{\ell})$ 作用在
$\cV_{\vec{\lambda}}^\dagger(\mathcal{F})$ 上。此外，若
$\theta(\vec{\ell})\ne 0$，则 $D(\vec{\ell})$ 作为扭转一阶微分算子作用在
$\cV_{\vec{\lambda}}^\dagger(\mathcal{F})$ 上。
\end{proposition}

\noindent\textit{证明.}\quad 取
\[
  \langle\tilde{\Psi}|\in\cV_{\vec{\lambda}}^\dagger(\mathcal{F}) .
\]
对任意元素 $X\in\mathfrak{g}$、$f\in\pi_*\cO_\mathcal{C}(*S)$ 与
$|\tilde{\Phi}\rangle\in\cH_{\vec{\lambda}}(\mathcal{B})$，由命题 4.7 与
(4.19) 我们有
\[
  \begin{aligned}
  \{D(\vec{\ell})\langle\tilde{\Psi}|\}(X\otimes f|\tilde{\Phi}\rangle)
  &= \theta(\vec{\ell})
     \bigl(\langle\tilde{\Psi}|X\otimes f|\tilde{\Phi}\rangle\bigr)
     -\langle\tilde{\Psi}|
     \bigl\{D(\vec{\ell})\bigl(X\otimes f|\tilde{\Phi}\rangle\bigr)\bigr\}\\
  &= 0 .
  \end{aligned}
\]
于是我们断定
\[
  D(\vec{\ell})\langle\Psi|\in\cV_{\vec{\lambda}}^\dagger(\mathcal{F}) .
\]
其余论断是定义 (4.18) 的容易推论。\hfill$\square$
% [Tutor Note: 第一行计算把 (4.19) 反用，并把两项分别化为零：(i) 因
%  $\langle\tilde\Psi|\in\cV^\dagger$ 是「在
%  $\hat{\mathfrak g}(\mathcal F)\cH$ 上取零」的泛函，函数
%  $\langle\tilde\Psi|X\otimes f|\tilde\Phi\rangle$ 恒等于 $0$，故
%  $\theta(\vec\ell)$ 作用其上得 $0$；(ii) 由命题 4.7，
%  $D(\vec\ell)(X\otimes f|\tilde\Phi\rangle)$ 仍属于
%  $\hat{\mathfrak g}(\mathcal F)\cH$，故 $\langle\tilde\Psi|$ 作用其上亦为
%  $0$。原文把这两步并成一行 "= 0."，是本块最容易卡住的地方。
%  命题 4.12 的意义也随之清楚：命题 4.7 只说 $D$ 在余真空层 $\cV$ 上良定义，
%  此处补上真空层 $\cV^\dagger$ 一侧；由于
%  $\cV^\dagger=\underline{\operatorname{Hom}}(\cV,\cO_\mathcal{B})$，
%  这本可由对偶自动得到，原文用 (4.19) 直接验证，是另一条等价的路。原文 p.82]

注意，对 (4.6) 中给出的双线性配对
$\langle\ |\ \rangle:\cV_{\vec{\lambda}}^\dagger(\mathcal{F})
\times\cV_{\vec{\lambda}}(\mathcal{F})\to\cO_\mathcal{B}$，我们有等式
\begin{equation}
  \{D(\vec{\ell})\langle\Psi|\}|\Phi\rangle
  +\langle\Psi|\{D(\vec{\ell})|\Phi\rangle\}
  = \theta(\vec{\ell})(\langle\Psi|\Phi\rangle) . \tag{4.23}
\end{equation}
